\section{Grundlagen}

Dieses Kapitel beschreibt die fachlichen und technischen Grundlagen, auf denen der in dieser Arbeit entwickelte Prototyp aufbaut. 
Zunächst werden grundlegende Konzepte der Zugriffskontrolle eingeordnet, insbesondere das Reference-Monitor-Konzept sowie attributbasierte Zugriffskontrollmodelle. 
Darauf aufbauend wird erläutert, wie Attribute in dynamischen Systemen als veränderliche Entscheidungsgrundlage verwendet werden können und welche Rolle textuelle Policy-Beschreibungen dabei einnehmen.

Für die Umsetzung im Linux-Kernel sind außerdem die Trennung zwischen Userspace und Kernelspace, Linux Security Modules sowie eBPF relevant. 
Diese Mechanismen bestimmen, an welchen Stellen Zugriffsentscheidungen technisch getroffen werden können und welche Einschränkungen bei der Implementierung im Kernel gelten. 
Ein besonderer Fokus liegt dabei auf eBPF-Maps als Schnittstelle zwischen Userspace und Kernelspace, gepinnten Maps für dauerhaft zugreifbare Zustände sowie den Vorgaben des eBPF-Verifiers.

Abschließend werden die für den Prototyp spezifischen Grundlagen betrachtet: die Verwendung von Aya als Rust-Framework, die Abbildung dynamischer Attribute und Policy-Generationen sowie die Überwachung bereits bestehender Zugriffe über File Descriptors. 
Damit schafft das Kapitel die Grundlage für die anschließende Anforderungsanalyse und die spätere Architektur des Systems.

\subsection{Zugriffskontrolle und Reference-Monitor-Konzept}

Zugriffskontrolle beschreibt Mechanismen, mit denen entschieden wird, ob ein Subjekt eine bestimmte Operation auf einer Ressource ausführen darf.
Ein Subjekt kann dabei beispielsweise ein Prozess oder ein Benutzer sein, während Ressourcen Dateien, Netzwerkendpunkte oder andere geschützte Objekte des Systems sein können.
Eine Zugriffsentscheidung basiert auf einer Menge von Regeln oder Policies, die festlegen, unter welchen Bedingungen eine Operation erlaubt oder verweigert wird.
Diese Sichtweise knüpft an grundlegende Begriffe der Informationsschutzmechanismen an, bei denen Akteure, geschützte Objekte und Berechtigungen als zentrale Elemente der Zugriffskontrolle betrachtet werden \citep{saltzer1975protection}.

Ein grundlegendes Konzept für die Durchsetzung von Zugriffskontrolle ist der Reference Monitor.
Der Reference Monitor ist eine abstrakte Sicherheitskomponente, die jeden Zugriff eines Subjekts auf ein geschütztes Objekt überprüft, bevor dieser Zugriff ausgeführt wird.
Damit ein Reference Monitor wirksam ist, muss er insbesondere drei Eigenschaften erfüllen: Er muss bei jedem relevanten Zugriff durchlaufen werden, gegen Manipulation geschützt sein und hinreichend klein beziehungsweise nachvollziehbar sein, damit seine Korrektheit überprüft werden kann \citep{anderson1972computer}.

Für Betriebssysteme ist dieses Konzept besonders relevant, da Zugriffe auf Dateien, Prozesse oder andere Systemressourcen letztlich über den Kernel vermittelt werden.
Wird Zugriffskontrolle erst außerhalb des Kernels umgesetzt, besteht die Gefahr, dass sicherheitsrelevante Operationen nicht vollständig erfasst oder durch alternative Zugriffspfade umgangen werden.
Eine Zugriffskontrolle im Kernel kann dagegen an zentralen Stellen ansetzen, an denen Operationen auf geschützte Ressourcen tatsächlich ausgeführt werden.
Dies steht im Zusammenhang mit dem Prinzip der vollständigen Vermittlung, nach dem jeder Zugriff auf ein geschütztes Objekt vor seiner Ausführung geprüft werden muss \citep{saltzer1975protection}.

Im Kontext dieser Arbeit bildet das Reference-Monitor-Konzept die Grundlage für die Frage, wie Attribute Stream-Based Access Control im Linux-Kernel umgesetzt werden kann.
Der Prototyp soll Zugriffe nicht nur anhand statischer Regeln bewerten, sondern dynamische Attribute in die Entscheidung einbeziehen.
Dadurch entsteht zusätzlich die Herausforderung, dass eine zunächst erlaubte Zugriffssituation später unzulässig werden kann, wenn sich entscheidungsrelevante Attribute ändern.
Das klassische Reference-Monitor-Konzept erklärt daher zunächst die unmittelbare Kontrolle eines Zugriffs, während die spätere Arbeit zusätzlich betrachtet, wie bestehende Zugriffssituationen nachträglich überwacht werden können.

\subsection{Attribute-Based Access Control}
Attribute-Based Access Control (ABAC) ist ein Zugriffskontrollmodell, bei dem eine Autorisierungsentscheidung nicht allein anhand einer festen Subjektidentität, wie bei ACL- oder identitätsbasierter Zugriffskontrolle, oder anhand einer zugewiesenen Rolle, wie bei Role-Based Access Control (RBAC), getroffen wird.
Stattdessen werden Attribute ausgewertet, die sich auf das Subjekt, die Ressource, die angeforderte Operation und gegebenenfalls auf die Umgebung beziehen.
Das National Institute of Standards and Technology (NIST) beschreibt ABAC als logische Zugriffskontrollmethode, bei der solche Attribute gegen Policies, Regeln oder Beziehungen geprüft werden, um zu entscheiden, ob eine Operation zulässig ist \citep{nistSP800162}.

Für diese Arbeit ist ABAC relevant, weil dadurch Zugriffsentscheidungen feiner modelliert werden können als bei rein rollenbasierten oder benutzerbasierten Ansätzen. 
Ein Subjekt kann beispielsweise über Attribute wie Position, Benutzerkennung oder Gruppenzugehörigkeit beschrieben werden. 
Eine Ressource kann Attribute wie Pfad, Typ oder Schutzstufe besitzen. 
Zusätzlich können Umwelt- oder Systemattribute, etwa Uhrzeit oder Systemzustand, in die Entscheidung einbezogen werden. 
Damit entsteht ein flexibles Modell, in dem Policies Bedingungen über mehrere Attributklassen ausdrücken können.

ABAC trennt damit die Beschreibung der Entscheidungsgrundlage von der konkreten technischen Durchsetzung. 
Die Policy beschreibt, unter welchen Bedingungen ein Zugriff erlaubt oder verweigert werden soll. 
Die Auswertung dieser Bedingungen und die Durchsetzung der Entscheidung können dagegen von unterschiedlichen Systemkomponenten übernommen werden. 
Diese Trennung ist für den Prototyp wichtig, weil die Policy-Verarbeitung im Userspace stattfinden kann, während die Zugriffsdurchsetzung später im Kernel betrachtet wird.

\subsection{Attribute Stream-Based Access Control}
Attribute Stream-Based Access Control (ASBAC) erweitert attributbasierte Zugriffskontrolle um fortlaufend veränderliche Entscheidungsgrundlagen und Autorisierungsentscheidungen.
Während eine klassische ABAC-Anfrage typischerweise zu einer einzelnen Entscheidung führt, können in ASBAC Attribute, Policies und weitere Entscheidungsgrundlagen als veränderliche Datenströme betrachtet werden.
Eine bestehende Autorisierung kann dadurch neu bewertet werden, sobald sich ein entscheidungsrelevanter Wert ändert.

Heutelbeck beschreibt ASBAC im Zusammenhang mit der Streaming Attribute Policy Language (SAPL) als Publish-Subscribe-Modell.
Ein Policy Enforcement Point stellt dazu eine Autorisierungsanfrage als Subscription an einen Policy Decision Point.
Der Policy Decision Point bezieht die für die Entscheidung benötigten Datenströme ein und kann dem Policy Enforcement Point aktualisierte Entscheidungen übermitteln, wenn sich Policies, Attribute oder weitere Entscheidungsgrundlagen ändern \citep{heutelbeck2021asbac}.

Die fortlaufende Neubewertung steht inhaltlich in der Nähe von Usage Control.
Unter Usage Control werden Zugriffskontrollmodelle verstanden, die eine Nutzung nicht nur zum Zeitpunkt der initialen Autorisierung betrachten, sondern auch Bedingungen während der laufenden Nutzung einbeziehen.
Usage Control erweitert klassische Zugriffskontrolle unter anderem um die Veränderbarkeit von Attributen und die Fortdauer von Entscheidungen während einer Nutzung.
Gouglidis et al. beschreiben diese Konzepte als \emph{attribute mutability} und \emph{continuity of decision} \citep{gouglidis2018usecon}.
Usage Control dient in dieser Arbeit damit als verwandtes theoretisches Konzept, ersetzt jedoch nicht die Definition von ASBAC.

Der in dieser Arbeit entwickelte Prototyp übernimmt nicht die vollständige Publish-Subscribe-Architektur von SAPL.
Insbesondere existiert kein zentraler PDP, der für jede bestehende Autorisierung einen individuellen Strom aktualisierter Entscheidungen an einen PEP übermittelt.
Stattdessen werden Policies und dynamische Attribute im Userspace verarbeitet und über eBPF-Maps als generationsbasierter Zustand bereitgestellt.
Neue Dateiöffnungen werden anhand dieses Zustands durch einen eBPF-LSM-basierten Kernelspace-Pfad bewertet.
Nach der erfolgreichen Aktivierung einer neuen Policy- oder Attributgeneration stößt der Prototyp ereignisbasiert eine Neubewertung bestehender Dateizugriffe an; zeitabhängige Policies lösen eine solche Neubewertung beim Erreichen einer entscheidungsrelevanten Zeitgrenze aus.

Ein geöffneter File Descriptor wird dabei als technische Repräsentation einer fortdauernden Nutzungssituation betrachtet, nicht als ASBAC-Subscription.
Der Prototyp registriert beim Öffnen keine fortdauernde Autorisierungsanfrage und verwaltet keinen individuellen Decision Stream.
Der Userspace rekonstruiert vielmehr nach einem Aktivierungs- oder Zeitereignis aus den dann sichtbaren Prozess- und File-Descriptor-Informationen eine Zugriffsanfrage und bewertet diese lokal erneut.
Der Prototyp realisiert somit keinen allgemeinen ASBAC-Decision-Stream, überträgt aber zentrale ASBAC-Eigenschaften auf einen begrenzten Zugriffskontrollmechanismus für Linux-Dateiöffnungen: veränderliche Entscheidungsgrundlagen, eine durch deren Änderung ausgelöste Neubewertung und die Reaktion auf eine geänderte Autorisierungsentscheidung.

\subsection{Policy-Beschreibungen für ASBAC}
Policies beschreiben die Bedingungen, unter denen eine Zugriffsentscheidung getroffen wird. 
In einem attributbasierten Modell beziehen sich diese Bedingungen auf Attribute von Subjekt, Ressource, Aktion und Umgebung. 
NIST beschreibt Policies in ABAC als Regeln oder Beziehungen, gegen die Attribute ausgewertet werden \citep{nistSP800162}. 
Für ASBAC bleibt diese Grundidee erhalten, sie wird jedoch um dynamisch veränderliche Attribute erweitert.

Eine Policy-Beschreibung muss dafür ausdrücken können, welche Aktion auf welche Ressource bezogen ist und welche Attributbedingungen erfüllt sein müssen. 
Im Prototyp kann dies beispielsweise bedeuten, dass das Öffnen einer bestimmten Datei nur erlaubt ist, wenn ein Subjektattribut einen bestimmten Wert besitzt. 
Die Policy-Datei bleibt dabei eine textuelle Beschreibung im Userspace. 
Für die Auswertung im Kernel muss diese Beschreibung später in eine kernelgeeignete Repräsentation überführt werden.

Etablierte Zugriffskontrollarchitekturen unterscheiden häufig zwischen logischen Rollen zur Entscheidung und zur Durchsetzung.
XACML beschreibt beispielsweise eine Architektur mit Policy Decision Point und Policy Enforcement Point sowie eine Policy-Sprache zur Beschreibung attributbasierter Entscheidungen \citep{oasisXacml30}.
Der Policy Decision Point (PDP) wertet die für eine Zugriffsanfrage anwendbaren Policies aus und erzeugt eine Autorisierungsentscheidung.
Der Policy Enforcement Point (PEP) schützt die Ressource, bildet eine Entscheidungsanfrage und setzt die erhaltene Entscheidung technisch durch.
PDP und PEP bezeichnen damit logische Rollen und müssen nicht notwendigerweise als getrennte Prozesse implementiert sein.

Diese Arbeit übernimmt weder XACML als Policy-Sprache noch dessen vollständiges Architekturmodell.
Im Prototyp sind Entscheidungs- und Durchsetzungsfunktion technisch eng gekoppelt.
Im Kernelspace erfassen der Einstieg und der Rückgabepfad des eBPF-LSM-Programms eine neue \texttt{file\_open}-Operation und setzen das Ergebnis über den LSM-Rückgabewert durch; sie erfüllen damit die PEP-Funktion.
Die dazwischenliegenden Stufen zur Auswahl und Auswertung statischer und dynamischer Policies sowie zum Combining erfüllen logisch eine integrierte PDP-Funktion.

Entsprechend rekonstruiert die Userspace-Komponente bestehende Zugriffssituationen aus Prozess- und File-Descriptor-Informationen, wertet sie lokal anhand der aktiven Policies und Attribute aus und versucht bei einer Deny-Entscheidung, den betroffenen File Descriptor zu schließen.
Auch diese Komponente verbindet somit eine Policy-Entscheidungsfunktion mit einer Durchsetzungsfunktion.
Die Bezeichnungen \emph{Kernelspace-PEP} und \emph{Userspace-PEP} werden im weiteren Verlauf beibehalten, um die jeweilige Durchsetzungsfunktion hervorzuheben.
Ein eigenständiger, von beiden PEPs getrennter PDP ist nicht Bestandteil des Prototyps.

\subsection{Userspace und Kernelspace in Linux}
Linux trennt zwischen Userspace und Kernelspace. 
Anwendungen laufen im Userspace und greifen auf Betriebssystemfunktionen über Systemaufrufe oder andere definierte Schnittstellen zu. 
Der Kernelspace enthält dagegen den Kernel und damit Code, der privilegiert ausgeführt wird und zentrale Ressourcen wie Prozesse, Speicher, Dateien und Geräte verwaltet. 
Die Kernel-Dokumentation beschreibt unterschiedliche Ausführungskontexte, darunter Code, der einem Prozess im Kernelspace zugeordnet ist, sowie Prozesse, die im Userspace laufen \citep{linuxKernelHacking}.

Diese Trennung ist für Sicherheitsmechanismen wesentlich. 
Komplexe Verarbeitung, Dateiparsing und Validierung sind im Userspace leichter zu entwickeln, zu testen und bei Fehlern weniger kritisch für die Stabilität des Gesamtsystems. 
Code im Kernelspace hat dagegen unmittelbaren Einfluss auf das Betriebssystem. 
Fehler in Kernelcode können nicht nur eine Anwendung, sondern das gesamte System beeinträchtigen. 
Daher ist es sinnvoll, umfangreiche Policy-Verarbeitung im Userspace zu belassen und den Kernelspace auf die sicherheitskritische Durchsetzung vorbereiteter Entscheidungen zu beschränken.

Für diese Arbeit ergibt sich daraus eine Architekturgrenze. 
Der Userspace liest Policies und Attribute ein, validiert sie und übersetzt sie in eine Form, die im Kernel effizient ausgewertet werden kann. 
Der Kernelspace übernimmt dagegen die Entscheidung an den relevanten Zugriffspunkten. 
Diese Aufteilung reduziert die Komplexität im Kernel, ohne die Durchsetzung vollständig aus dem sicherheitsrelevanten Ausführungspfad herauszulösen.

\subsection{Linux Security Modules und LSM-Hooks}
Linux Security Modules (LSM) stellen ein Framework bereit, über das Sicherheitsmodule an sicherheitsrelevanten Stellen des Kernels Entscheidungen treffen können. 
Die Kernel-Dokumentation beschreibt LSM als Schnittstelle für Sicherheitsmodule wie SELinux, Smack, AppArmor oder TOMOYO \citep{kernelorgLSM}. 
Aus technischer Sicht werden dafür Hooks verwendet, die an bestimmten Punkten im Kernel-Ausführungspfad aufgerufen werden.

LSM-Hooks sind für Zugriffskontrolle besonders relevant, weil sie Operationen an Stellen erfassen können, an denen der Kernel ohnehin über die Zulässigkeit einer Operation entscheidet. 
Wright et al. beschreiben das LSM-Framework als allgemeine Sicherheitsunterstützung im Linux-Kernel, bei der Hooks Zugriff auf Kernelobjekte vermitteln und Rückgabewerte über Erfolg oder Fehler der Operation entscheiden können \citep{wright2002linux}. 
Damit eignen sich LSM-Hooks grundsätzlich als technische Grundlage für einen Reference-Monitor-ähnlichen Durchsetzungspunkt.

Für den Prototyp ist insbesondere der Hook \texttt{file\_open} relevant. 
Er erlaubt eine Kontrolle beim Öffnen einer Datei und passt damit zur funktionalen Einschränkung dieser Arbeit auf Dateiöffnungen. 
LSM-Hooks bewerten jedoch jeweils konkrete Kerneloperationen. 
Ein beim Öffnen erlaubter Zugriff wird nicht automatisch erneut durch denselben Hook geprüft, wenn sich später Attribute ändern. 
Diese Eigenschaft erklärt, warum zusätzlich eine Überwachung bestehender Zugriffssituationen betrachtet werden muss.

\subsection{eBPF und eBPF-LSM}
eBPF ist ein Mechanismus, mit dem Programme in einer eingeschränkten Ausführungsumgebung im Linux-Kernel ausgeführt werden können. 
Die Programme werden vor der Ausführung durch den Kernel geprüft und können an verschiedene Ereignisse oder Hook-Typen gebunden werden. 
Historisch steht BPF für Berkeley Packet Filter. 
Die heute im Linux-Kernel verwendete erweiterte Form wird meist als eBPF bezeichnet; die Kernel-Dokumentation verwendet für das moderne Subsystem jedoch häufig weiterhin die Bezeichnung BPF. 
In dieser Arbeit wird deshalb eBPF verwendet, wenn der moderne Mechanismus für Programme, Maps und Verifier gemeint ist. 
Der Begriff BPF wird nur dort aufgegriffen, wo er in offiziellen Bezeichnungen wie \texttt{bpf()}-Systemaufruf, BPF-Dateisystem oder in der Kernel-Dokumentation selbst verwendet wird.
Die Kernel-Dokumentation beschreibt BPF als im Kernel verwendetes Instruktionsformat und betont, dass BPF-Programme auf definierte Hilfsfunktionen und Kernelobjekte wie Maps beschränkt sind \citep{kernelBpfDesignQA}.

eBPF-LSM verbindet eBPF mit dem LSM-Framework. 
Die Kernel-Dokumentation beschreibt LSM-BPF-Programme als Möglichkeit, LSM-Hooks zur Laufzeit durch privilegierte Benutzer zu instrumentieren, um systemweite Mandatory-Access-Control- oder Audit-Policies mit eBPF umzusetzen \citep{kernelBpfLsm}. 
Damit kann ein eBPF-Programm an einem LSM-Hook ausgeführt werden und dort eine Entscheidung über die angeforderte Operation zurückgeben.

Für das Laden und die typgerechte Anbindung solcher Programme ist außerdem das \emph{BPF Type Format} (BTF) relevant.
BTF ist ein Metadatenformat, das unter anderem Typ- und Funktionsinformationen zu Kernel- und eBPF-Strukturen beschreibt \citep{kernelBtf}.
Ein Loader kann die vom Zielkernel unter \path{/sys/kernel/btf/vmlinux} bereitgestellten Informationen verwenden, um den vorgesehenen LSM-Hook und die Typen seiner Argumente aufzulösen.
BTF ist damit keine Policy- oder Attributquelle, sondern eine technische Voraussetzung für das Laden des in dieser Arbeit eingesetzten eBPF-LSM-Programms.

In den Grundlagen wird eBPF noch nicht als endgültige Lösung begründet, sondern als technischer Mechanismus eingeordnet. 
Für diese Arbeit ist eBPF interessant, weil es eine Ausführung im Kernel ermöglicht, ohne ein klassisches Kernelmodul zu entwickeln. 
Gleichzeitig ist eBPF durch Verifier, eingeschränkte Sprachelemente, begrenzten Stack, fehlende allgemeine Heap-Allokation und definierte Hilfsfunktionen deutlich stärker beschränkt als allgemeiner Kernelcode. 
Diese Einschränkungen müssen in der späteren Konzeption gegen die Vorteile abgewogen werden.

\subsection{eBPF-Maps und gepinnte Maps}
eBPF-Maps sind Datenstrukturen im Kernel, die von eBPF-Programmen und Userspace-Anwendungen gemeinsam genutzt werden können. 
Die Kernel-Dokumentation beschreibt Maps als generischen Speicher für unterschiedliche Datentypen, der unter anderem den Datenaustausch zwischen Kernelspace und Userspace ermöglicht \citep{kernelBpfMaps}. 
Über den Systemaufruf \texttt{bpf()} können Maps erzeugt sowie Einträge gelesen, aktualisiert oder gelöscht werden \citep{manBpf}.

Für Zugriffskontrolle sind Maps relevant, weil eBPF-Programme nicht beliebige komplexe Userspace-Datenstrukturen direkt verwenden können. 
Policies und Attribute müssen daher in eine feste, kernelgeeignete Struktur gebracht werden. 
Hash-Maps oder Arrays können beispielsweise verwendet werden, um Policy-Einträge, Attributwerte, Generationen oder Statusinformationen bereitzustellen. 
Der Kernel kann diese Werte während einer Zugriffskontrolle auslesen, während der Userspace sie bei Policy- oder Attributänderungen aktualisiert.

Gepinnte Maps erweitern diese Idee um einen stabilen Zugriffspfad im BPF-Dateisystem. 
Pinning erlaubt es, eBPF-Objekte wie Maps über einen Pfad zugänglich zu machen, sodass andere Prozesse sie später wieder öffnen können \citep{ebpfPinning}. 
Dies ist für den Prototyp wichtig, weil nicht nur der Loader, sondern auch ein Userspace-PEP oder Administrationswerkzeug den aktuellen Zustand lesen können muss.
Gepinnte Maps bilden damit eine zentrale Schnittstelle zwischen den getrennten Userspace-Komponenten und dem Kernelzustand.

\subsection{eBPF-Verifier, Stack und Programmkomplexität}
Bevor ein eBPF-Programm im Kernel ausgeführt werden darf, prüft der eBPF-Verifier dessen Sicherheit. 
Die Kernel-Dokumentation beschreibt die Prüfung als mehrstufigen Prozess, bei dem unter anderem Kontrollfluss, Registerzustände, Stackzugriffe und Speicherzugriffe analysiert werden \citep{kernelBpfVerifier}. 
Ein Programm wird nur geladen, wenn der Verifier zeigen kann, dass es die Sicherheitsregeln erfüllt.

Diese Prüfung hat direkte Auswirkungen auf die Gestaltung von eBPF-Programmen. 
Zugriffe auf Map-Werte, Stackbereiche oder Kontextdaten müssen für den Verifier nachvollziehbar sein. 
Der Verifier verfolgt unter anderem, ob Register initialisiert sind, welche Zeigerarten verwendet werden und ob Speicherzugriffe innerhalb bekannter Grenzen liegen. 
Unklare oder zu komplexe Zugriffsmuster können dazu führen, dass ein Programm abgelehnt wird, auch wenn die fachliche Logik korrekt erscheint.
Darüber hinaus muss der Kontrollfluss eines eBPF-Programms für den Verifier analysierbar bleiben.
Schleifen dürfen nicht unbeschränkt sein, Rekursion ist nicht als allgemeines Programmiermittel verfügbar, und Speicherzugriffe müssen durch bekannte Grenzen abgesichert sein.
Auch der Zugriff auf Kernelinformationen erfolgt nicht beliebig, sondern über den Hook-Kontext, Maps oder für den jeweiligen Programmtyp zugelassene Hilfsfunktionen.

Zusätzlich bestehen praktische Begrenzungen bei Stack und Programmkomplexität. 
Die BPF-Dokumentation nennt für Programme ein Stacklimit von 512 Byte und beschreibt interne Grenzen des Verifiers, darunter die Anzahl der während der Analyse betrachteten Instruktionen und Zustände \citep{kernelBpfDesignQA}. 
Die offizielle Aya-Dokumentation beschreibt eBPF als eingeschränkte Laufzeitumgebung und weist darauf hin, dass eBPF-Programme nur 512 Byte Stack besitzen; bei Programmen, die Tail Calls verwenden, kann der verfügbare Stack auf 256 Byte begrenzt sein \citep{ayaBook}. 
Der Stack eignet sich damit nur für kleine lokale Hilfsdaten, etwa Schlüsselstrukturen für Map-Zugriffe oder wenige temporäre Werte. 
Größere lokale Puffer, dynamisch wachsende Listen oder Zeichenketten können dort nicht sinnvoll abgelegt werden. 
Ein allgemeiner Heap, wie er in Userspace-Programmen für dynamische Datenstrukturen wie Listen, Strings oder Objekte genutzt wird, steht eBPF-Programmen nicht zur Verfügung. 
Veränderliche oder größere Zustände müssen daher typischerweise in eBPF-Maps ausgelagert werden.
Für Rust-basierte eBPF-Programme ergeben sich daraus weitere Einschränkungen.
Die Standardbibliothek kann im eBPF-Programm nicht verwendet werden; stattdessen steht nur \texttt{core} zur Verfügung.
Funktionen und Traits, die auf \texttt{core::fmt} beruhen, etwa \texttt{Display} oder \texttt{Debug}, sind nicht nutzbar.
Da kein Heap verfügbar ist, können auch \texttt{alloc} und darauf aufbauende Collections nicht eingesetzt werden.
Außerdem darf ein eBPF-Programm nicht auf Panics angewiesen sein, weil die eBPF-Laufzeit weder Stack-Unwinding noch eine Abort-Instruktion unterstützt.
Ein eBPF-Programm besitzt zudem keine gewöhnliche \texttt{main}-Funktion, sondern wird über einen programmspezifischen Einstiegspunkt an einen Hook gebunden.
Da eBPF-Programme häufig direkt Kernel-Speicher oder Hook-Kontextdaten lesen, ist außerdem ein relevanter Teil des Codes als \texttt{unsafe} zu behandeln \citep{ayaBook}.
Für die Arbeit bedeutet dies, dass komplexe Policy-Logik möglichst im Userspace vorbereitet werden sollte, während das eBPF-Programm nur eine kompakte, vorher strukturierte Repräsentation auswertet.

\subsection{Aya als Rust-Framework für eBPF}
Aya ist ein Rust-Framework zur Entwicklung von eBPF-Anwendungen. 
Die offizielle Aya-Dokumentation beschreibt Aya als Möglichkeit, eBPF-Programme mit Rust zu entwickeln und dabei Cargo, Rust-Strukturen sowie gemeinsame Codeanteile zwischen Userspace und eBPF-Programm zu nutzen \citep{ayaBook}. 
Neben Rust-basierten Frameworks existieren etablierte eBPF-Werkzeuge im C- und Skriptsprachenumfeld.
libbpf ist eine C-Bibliothek für die Entwicklung und das Laden von BPF- beziehungsweise eBPF-Programmen, während BCC (BPF Compiler Collection) ein Toolkit für BPF-basierte Analyse-, Tracing- und Monitoring-Werkzeuge ist \citep{libbpf,bcc}.
Aya stellt dazu eine Rust-basierte Alternative dar.

Für den Prototyp ist Aya vor allem deshalb relevant, weil Userspace-Komponenten und eBPF-Programm in derselben Programmiersprache entwickelt werden können. 
Gemeinsame Datentypen können in einer gemeinsamen Bibliothek definiert werden, wodurch sich die Gefahr verringert, dass Userspace und Kernelspace unterschiedliche Annahmen über die Struktur von Map-Schlüsseln oder Map-Werten treffen. 
Gerade bei Policies und Attributen ist diese Konsistenz wichtig, weil dieselben binären Strukturen auf beiden Seiten der Map-Schnittstelle interpretiert werden.

Gleichzeitig hebt Aya die Einschränkungen von eBPF nicht auf. 
Auch mit Rust gelten die Vorgaben des eBPF-Verifiers, die Stackbegrenzung, der Verzicht auf Heap-Allokation im eBPF-Programm und die Beschränkung auf unterstützte Hilfsfunktionen. 
Aya ist daher als Entwicklungswerkzeug zu verstehen, nicht als fachliches Zugriffskontrollmodell. 
Die Entscheidung für Aya betrifft die spätere Implementierung, während die fachlichen Anforderungen weiterhin aus ASBAC und der Kernel-Durchsetzung abgeleitet werden.

\subsection{Dynamische Attribute und Policy-Generationen}
Dynamische Attribute sind Attributwerte, die sich zur Laufzeit ändern können. 
Im Unterschied zu statischen Eigenschaften, etwa einem festen Dateipfad, können dynamische Attribute zeitabhängig oder durch externe Ereignisse beeinflusst sein. 
Beispiele sind Uhrzeit, Systemzustand oder vom Nutzer frei definierte Subjektattribute. 
Die Idee veränderlicher Attribute knüpft an Usage-Control-Modelle an, in denen Attributmutabilität als ausdrückliches Merkmal betrachtet wird \citep{gouglidis2018usecon}.

Für den Prototyp bedeutet dies, dass Attribute nicht nur beim Start einmalig geladen werden dürfen. 
Änderungen müssen vom Userspace erkannt, validiert und in die kernelgeeignete Repräsentation übertragen werden. 
Da eBPF-Programme zur Laufzeit auf Maps zugreifen können, bieten Maps eine geeignete Schnittstelle, um aktuelle Attributwerte für Zugriffsentscheidungen bereitzustellen \citep{kernelBpfMaps}. 
Die Attributnamen und Werte müssen dabei im Userspace so verarbeitet werden, dass das eBPF-Programm mit festen Schlüssel- und Wertstrukturen arbeiten kann.

Policy-Generationen beschreiben zusammengehörige gültige Stände von Policies. 
Dieser Begriff ist für den Prototyp wichtig, weil Policy-Dateien im Userspace geparst und validiert werden, bevor sie für Zugriffsentscheidungen verwendet werden. 
Eine neue Generation sollte erst dann aktiv werden, wenn alle dazugehörigen Map-Einträge vollständig erzeugt und konsistent bereitgestellt wurden. 
Dadurch wird vermieden, dass der Kernel während einer Aktualisierung auf einen unvollständigen oder widersprüchlichen Policy-Zustand zugreift.

\subsection{Überwachung bestehender Zugriffe und File Descriptors}
Ein Dateizugriff in Linux wird häufig über einen File Descriptor repräsentiert. 
Der Systemaufruf \texttt{open()} liefert bei Erfolg einen nichtnegativen File Descriptor zurück, der in nachfolgenden Systemaufrufen wie \texttt{read()}, \texttt{write()} oder \texttt{lseek()} als Referenz auf die geöffnete Datei verwendet wird \citep{manOpen}. 
Damit entsteht nach dem Öffnen eine bestehende Zugriffssituation, die unabhängig vom ursprünglichen Pfad weiter über den File Descriptor genutzt werden kann.

Für eine Zugriffskontrolle über \texttt{file\_open} ist das wichtig, weil der LSM-Hook beim Öffnen der Datei greift. 
Wenn sich danach entscheidungsrelevante Attribute ändern, wird derselbe bereits geöffnete File Descriptor nicht automatisch erneut durch \texttt{file\_open} bewertet. 
Eine ASBAC-orientierte Zugriffskontrolle muss deshalb zusätzlich betrachten, wie bestehende File Descriptors überwacht und bei geänderter Entscheidungsgrundlage behandelt werden können.

Linux stellt über \texttt{/proc/<pid>/fd} Informationen über die geöffneten File Descriptors eines Prozesses bereit. 
Die man-pages beschreiben die Einträge in diesem Verzeichnis als symbolische Links auf die geöffneten Dateien oder andere referenzierte Objekte; der Zugriff darauf unterliegt eigenen Berechtigungsprüfungen \citep{manProcPidFd}. 
Für den Prototyp kann ein Userspace-PEP diese Informationen nutzen, um zu erkennen, welcher Prozess über welchen File Descriptor auf welche Datei zugreift.
Wird eine bestehende Zugriffssituation durch neue Attribute oder Policies unzulässig, kann der Userspace-PEP gezielt den betroffenen Zugriff behandeln, ohne notwendigerweise alle Zugriffe desselben Prozesses zu beenden.

\newpage
